\input{../tex-inputs/latex-korrekturansicht-vorspann}

\section[Theodor Herzl an Arthur Schnitzler, 14. 8. 1885]{L03821 Theodor Herzl an Arthur Schnitzler, 14. 8. 1885}
\nopagebreak\mylabel{L03821v}
\rehead{ }\normalsize\beginnumbering\briefempfaengerindex{Schnitzler, Arthur@\textsc{Schnitzler, Arthur}!zzzHerzl, Theodor@\emph{von Theodor Herzl}!1885-08-141@{14. 8. 1885}|(be}
\toendnotes[C]{\smallbreak\pagebreak[2]}
\correspDesc{Versand  durch Theodor Herzl am 14. 8. 1885 in Schörfling
\newline{}Erhalt  durch Arthur Schnitzler im Zeitraum [14. 8. 1885
                  – 15. 8. 1885?] in Innsbruck}\toendnotes[C]{\smallbreak}
\Standort{CUL, Schnitzler, B 39.}
\physDesc{Brief, 1 Blatt, 1 Seite, 545 Zeichen
\newline{}Handschrift: schwarze Tinte, lateinische Kurrent
\newline{}Ordnung: mit Bleistift von unbekannter Hand nummeriert: »2« }
\buchAbdrucke{\weitereDrucke{Theodor Herzl: \emph{Briefe und autobiographische Notizen 1866–1895}. Bearbeitet von Johannes Wachten in Zusammenarbeit mit Chaya Harel, Daisy Tycho und Manfred Winkler. Berlin, Frankfurt am Main, Wien: \emph{Propyläen} 1983, S. 199 (Briefe und Tagebücher. Herausgegeben von Alex Bein, Hermann Greive, Moshe Schaerf, Julius H. Schoeps und Johannes Wachten, 1).} }\toendnotes[C]{\smallbreak}
\pstart
           \raggedleft{}{\pb}\textcolor{pink}{Hotel Kammer}\oindex{Hotel Kammer@\textbf{Hotel Kammer}, \emph{Hotel}|pw}{}\ledrightnote{\textcolor{pink}{Hotel Kammer}}, 14/8 885\pend
           
\pstart{}Theuerster Doctor!\pend\vspace{0.5em}
\pstart
           Ihre \label{K_L03821-1v}\edtext{Depesche, Ihren lieben Brief}{\lemma{\textnormal{\emph{Depesche, … Brief}}}\Cendnote{\textnormal{Depesche: ein Eilschreiben, vermutlich ein
                  Telegramm. Dieses ist, ebenso wie der erwähnte Brief, nicht überliefert. Laut \emph{\textcolor{green}{Tagebuch}\pwindex{Schnitzler, Arthur 15. 5. 1862 Wien – 21. 10. 1931 ebd.@\textsc{Schnitzler, Arthur} (15. 5. 1862 Wien – 21. 10. 1931 ebd.), \emph{Schriftsteller, Mediziner}!Tagebuch@\strich\emph{Tagebuch}|pwk}} hatte \textcolor{blue}{Schnitzler}{ }\textcolor{blue}{Herzl}\pwindex{Herzl, Theodor 2.\,5.\,1860 Budapest – 3.\,7.\,1904 Edlach@\textsc{Herzl, Theodor} (2.\,5.\,1860 Budapest – 3.\,7.\,1904 Edlach), \emph{Schriftsteller, Journalist}|pwk} am 8. 8. 1885 in \textcolor{pink}{Kammer}\oindex{Kammer@\textbf{Kammer}|pwk} getroffen. \textcolor{blue}{Schnitzler} dürfte 
                  \textcolor{blue}{Herzl}\pwindex{Herzl, Theodor 2.\,5.\,1860 Budapest – 3.\,7.\,1904 Edlach@\textsc{Herzl, Theodor} (2.\,5.\,1860 Budapest – 3.\,7.\,1904 Edlach), \emph{Schriftsteller, Journalist}|pwk} eingeladen haben, am 14. 8. 1885 gemeinsam abzureisen. \textcolor{blue}{Schnitzler} reiste von \textcolor{pink}{Ischl}\oindex{Bad Ischl@\textbf{Bad Ischl}|pwk}
                  nach \textcolor{pink}{Innsbruck}\oindex{Innsbruck@\textbf{Innsbruck}, \emph{Verwaltungsgebiet}|pwk}. Aber \textcolor{blue}{Herzl}\pwindex{Herzl, Theodor 2.\,5.\,1860 Budapest – 3.\,7.\,1904 Edlach@\textsc{Herzl, Theodor} (2.\,5.\,1860 Budapest – 3.\,7.\,1904 Edlach), \emph{Schriftsteller, Journalist}|pwk}, der in \textcolor{pink}{Kammer}\oindex{Kammer@\textbf{Kammer}|pwk}
                  seine \textcolor{blue}{Eltern}\pwindex{Herzl, Jeanette 28.\,7.\,1836 Budapest – 20.\,2.\,1911 Wien@\textsc{Herzl, Jeanette} (28.\,7.\,1836 Budapest – 20.\,2.\,1911 Wien)|pwkv}\pwindex{Herzl, Jakob 14.\,3.\,1837 Zemun – 9.\,6.\,1902 Wien@\textsc{Herzl, Jakob} (14.\,3.\,1837 Zemun – 9.\,6.\,1902 Wien), \emph{Bankdirektor, Großkaufmann}|pwkv}
                  besuchte, trat seine Reise über \textcolor{pink}{München}\oindex{München@\textbf{München}|pwk} erst
                  am Folgetag an.}}}\label{K_L03821-1} habe ich mit Vergnügen und Bedauern erhalten.
               Letzteres hat seinen Grund darin, dass ich erst morgen mich auf die
                  \label{K_L03821-2v}\edtext{\textcolor{pink}{holländischen}\oindex{Niederlande@\textbf{Niederlande}|pw}{}\ledrightnote{\textcolor{pink}{Niederlande}} Strümpfe}{\lemma{\textnormal{\emph{holländischen Strümpfe}}}\Cendnote{\textnormal{\textcolor{blue}{Herzl}\pwindex{Herzl, Theodor 2.\,5.\,1860 Budapest – 3.\,7.\,1904 Edlach@\textsc{Herzl, Theodor} (2.\,5.\,1860 Budapest – 3.\,7.\,1904 Edlach), \emph{Schriftsteller, Journalist}|pwk} brach am 15. 8. 1885 zu
                  einer Studien- und journalistischen Reise nach \textcolor{pink}{Belgien}\oindex{Belgien@\textbf{Belgien}|pwk} und \textcolor{pink}{Holland}\oindex{Niederlande@\textbf{Niederlande}|pwk} auf, von der er
                  am 11. 9. 1885 zurückkehrte.}}}\label{K_L03821-2} mache. Ich konnte meinen guten
                  \textcolor{blue}{Alten}\pwindex{Herzl, Jeanette 28.\,7.\,1836 Budapest – 20.\,2.\,1911 Wien@\textsc{Herzl, Jeanette} (28.\,7.\,1836 Budapest – 20.\,2.\,1911 Wien)|pw}\pwindex{Herzl, Jakob 14.\,3.\,1837 Zemun – 9.\,6.\,1902 Wien@\textsc{Herzl, Jakob} (14.\,3.\,1837 Zemun – 9.\,6.\,1902 Wien), \emph{Bankdirektor, Großkaufmann}|pw}{}\ledrightnote{\textcolor{blue}{Jeanette Herzl}{\newline}\textcolor{blue}{Jakob Herzl}} nicht diesen Tag stehlen,
               obgleich im Tagediebstahl bereits Einiges geleistet habe.\pend
           
\pstart
           Ihre Depesche konnte ich nicht erwidern, weil ich Ihre Adresse nicht hatte. –
               Nochmals besten Dank!\pend
           
\pstart
           Ich wünsche Ihnen eine sonnige, vergnügte Fahrt!\pend
           
\pstart
           Glauben Sie an die freundschaftlichen Gefühle Ihres aufrichtig ergebenen{\\[\baselineskip]}\spacefill\mbox{Herzl}\pend
           \leftskip=0em{}\selectlanguage{ngerman}\endnumbering\briefempfaengerindex{Schnitzler, Arthur@\textsc{Schnitzler, Arthur}!zzzHerzl, Theodor@\emph{von Theodor Herzl}!1885-08-141@{14. 8. 1885}|)be}\mylabel{L03821h}  \newcommand{\dateiname}{L03821}\newcommand{\titel}{Theodor Herzl an Arthur Schnitzler, 14. 8. 1885}\newcommand{\editorInnen}{Selma Jahnke und Martin Anton Müller}\input{../tex-inputs/latex-korrekturansicht-abspann}
